\documentclass[11pt,a4paper]{article}
\usepackage[T1]{fontenc}
\usepackage{lmodern,amsmath,amssymb,amsthm,mathtools,microtype}
\usepackage[margin=27mm]{geometry}
\usepackage[colorlinks=true]{hyperref}
\setlength{\emergencystretch}{2em}
\newtheorem{proposition}{Proposition}[section]
\newtheorem{corollary}[proposition]{Corollary}
\newcommand{\Q}{\mathbb Q}
\title{A quadratic trace class and the fresh-contact trace}
\author{}\date{17 November 2026}
\begin{document}\maketitle
\begin{abstract}
A quadratic trace class of coefficient height three preserves the
negative-odd-pole test and all its prepared local square conditions.
We also determine the trace residue at every fresh comparison contact:
it is one. Neither assertion proves simultaneous rational realization
on actual detector tuples or a smaller complete defining-variable bound.
\end{abstract}

\section{A smaller trace class}
Let $S$ contain two. Let $w,V\in\Q^\times$ and assume that every
odd valuation of $w$ outside $S$ is positive. Put
\[
 b=4w/V^2,\qquad C=V^2+w(2-V).
\]
For a fixed positive rational $S$-unit $t$ with $0<t<8$, put
\[
 D_t(V)=V^2+t(V-1)(V-2).
\]
\begin{proposition}
If $D_t(V)C$ is a rational square, then $b$ has no negative odd
valuation outside $S$. The polynomial $D_t$ is strictly positive on
the whole real line, and
\[
 D_t(0)=2t,\qquad D_t(1)=1,\qquad D_t(2)=4.
\]
Consequently this class has the same fixed-place and selected-contact
local square preparation as the trace-conic criterion. It introduces
no additional free coordinate or exceptional prime.
\end{proposition}
\begin{proof}
At a hypothetical negative odd pole, write $\nu=v_p(w)>0$ odd and
$m=v_p(V)>\nu/2$. Since $p\ne2$, the term $w(2-V)$ has
valuation $\nu$, strictly less than $2m$, so $v_p(C)=\nu$.
Because $t$ is an $S$-unit, $D_t(V)\equiv2t\not\equiv0\pmod p$.
The product still has odd valuation and cannot be a square.

The leading coefficient of $D_t$ is $1+t>0$ and its discriminant
is $t(t-8)<0$. This proves strict real positivity. The displayed
three evaluations follow directly. Near the fixed reference trace
$V=2$, the class is near the nonzero rational square four. At a
selected contact with $V\equiv1$ or $2\pmod p$, it has residue
one or four. The usual local square neighborhoods give the assertions,
including arbitrary fixed precision at two. All denominators are
already those of $V$ or fixed $S$-units. The finite-algebra square
extension and guard norms therefore have the same coordinate count.
\end{proof}

For $t=1$ this gives
\[
 \boxed{D_1(V)=2V^2-3V+2.}
\]
The coefficient height is three. In the class of integer polynomials
of degree at most two, with nonzero constant coefficient, positive on
$[0,2]$, and nonzero square values at both one and two, every example
of coefficient height at most two is either the constant one or the
square $(V+1)^2$. Thus height three is minimal for a nonsquare
polynomial in this specified class. This finite assertion is verified
by the accompanying exact checker.

\section{The trace at an actual fresh comparison contact}
Let $K=\Q(\sqrt2)$, and retain the general or conditional source
with its reference point $P$ and product $W$. At a fresh split
comparison prime $p$, the curve has good reduction, $P$ has finite
reduction at both places, and $v_p(w)=1$. Each comparison denominator
has odd valuation, with its other normalizers units. The tests include
$f(H)$ and $\overline{f(H-P)}$, where $f=1/x$.

\begin{proposition}
At every fresh selected comparison prime in the retained completeness
construction,
\[
                 \operatorname{Tr}_{K/\Q}(W^{-1})\equiv1\pmod p.
\]
The alternative residue two does not occur there.
\end{proposition}
\begin{proof}
The successful slot has $W\equiv1\pmod p$. Every factor in the
other slot has nonpositive valuation: its second summand has odd
valuation and cannot cancel its constant term. If that product were
a unit, every factor would be a unit and every field test would have
positive valuation at both places. In particular $f(H)$ and
$f(H-P)$ would have positive valuation at the same place.
At good reduction this means that both $H$ and $H-P$ reduce to the
identity. Their difference $P$ would also reduce to the identity,
contrary to its finite reduction. Thus the other product has negative
valuation and its reciprocal vanishes modulo $p$. Adding the two
reciprocal slots gives one.

The extra rational factors in the general detector have nonpositive
valuations as well, so they cannot undo this argument. The selected
large primes avoid the fixed denominators of $P$ and the bad-reduction
set; no additional exclusion is needed.
\end{proof}

The conic, the quadratic-class elliptic tests and adaptive trace
classes give separate arithmetic possibilities. Their simultaneous
rational realization at the actual prescribed detector traces remains
unproved.
\end{document}
